\section{Pendahuluan}

Kriptografi merupakan ilmu dan seni untuk menjaga kerahasiaan pesan dengan cara menyandikan pesan tersebut sehingga hanya pihak yang berwenang yang dapat membacanya \cite{katz2020}. Secara etimologi, kata ``kriptografi'' berasal dari bahasa Yunani: \textit{kryptos} (tersembunyi) dan \textit{graphein} (menulis). Dengan demikian, kriptografi dapat diartikan sebagai ``tulisan tersembunyi''.

Dalam era digital saat ini, kriptografi menjadi fondasi utama keamanan informasi \cite{stallings2017}. Setiap aktivitas komunikasi daring---mulai dari pesan instan, transaksi perbankan, hingga penyimpanan data---mengandalkan teknik kriptografi untuk melindungi data dari pihak yang tidak berwenang.

Buku ajar ini dirancang untuk mahasiswa Program Studi Matematika jenjang S1 yang mempelajari kriptografi sebagai mata kuliah wajib atau pilihan. Materi disusun secara sistematis mulai dari konsep dasar, matematika pendukung, algoritma klasik dan modern, hingga implementasi dalam berbagai bahasa pemrograman (MATLAB, C, C++, dan Pascal).
